% Search from the source directory or the repository root.
\makeatletter
\def\input@path{{../../preamble/}{book2026-2027/preamble/}}
\makeatother
\input{preamble.tex}
\usepackage[spanish,es-nodecimaldot]{babel}
\usepackage{lmodern}

\geometry{letterpaper,margin=0.25in}
\hypersetup{pdfprintscaling=None}
\pagestyle{empty}
\linespread{1}
\setlength{\parskip}{0pt}
\setlength{\tabcolsep}{2.5pt}
\setlength{\arrayrulewidth}{0.4pt}
\setlength{\extrarowheight}{0.7pt}
\renewcommand{\arraystretch}{1}
\newcolumntype{R}{>{\centering\arraybackslash}p{38pt}}

\definecolor{gridgray}{gray}{0.55}
\arrayrulecolor{gridgray}

\newcommand{\writingrule}{%
  \leaders\hrule height 0.1pt depth 0.3pt\hfill\kern0pt}
\newcommand{\field}[1]{%
  \rule[-2pt]{0pt}{15pt}\textbf{#1}\enspace\writingrule}

\newcommand{\category}[1]{%
  \multicolumn{5}{|l|}{\cellcolor{black!9}\textbf{#1}}\\\hline}
\newcommand{\criterion}[1]{#1 & & & & \\\hline}
\newcommand{\projectheading}[1]{{\scriptsize\bfseries Proyecto\newline #1}}

\newcommand{\responseblock}[1]{%
  \fcolorbox{gridgray}{white}{%
    \begin{minipage}[t]{\dimexpr0.49\linewidth-2\fboxsep-2\fboxrule\relax}
      \textbf{Proyecto #1}\par
      \field{Problemática}\par
      \field{Fortaleza}\par
      \field{Mejora concreta}\par
      \field{Solución alternativa}\par
      \field{Cómo o por qué}\par
    \end{minipage}}%
}

\begin{document}

{\Large\bfseries Coevaluación de la feria\par}
\vspace{3pt}
\begin{minipage}[t]{0.72\linewidth}
  \field{Grupo evaluador}
\end{minipage}\hfill
\begin{minipage}[t]{0.25\linewidth}
  \field{Curso}
\end{minipage}\par

\textbf{Nombre de cada proyecto o grupo evaluado}\par
\begin{minipage}[t]{0.235\linewidth}\field{1}\end{minipage}\hfill
\begin{minipage}[t]{0.235\linewidth}\field{2}\end{minipage}\hfill
\begin{minipage}[t]{0.235\linewidth}\field{3}\end{minipage}\hfill
\begin{minipage}[t]{0.235\linewidth}\field{4}\end{minipage}\par

\vspace{5pt}
\fontsize{10bp}{11.3pt}\selectfont
Escribe una letra por criterio en la columna del proyecto que estés observando.\par
\textbf{E} = Se evidencia \enspace
\textbf{P} = Parcialmente \enspace
\textbf{A} = Aún no se evidencia \enspace
\textbf{N} = No se observa\par
\vspace{3pt}

\renewcommand{\arraystretch}{1.1}

\begin{tabularx}{\linewidth}{|>{\raggedright\arraybackslash}X|R|R|R|R|}
  \hline
  \rowcolor{black!14}
  \textbf{Criterios} & \projectheading{1} & \projectheading{2} &
  \projectheading{3} & \projectheading{4} \\\hline
  \category{1. Comprensión de la problemática}
  \criterion{El grupo explica claramente la problemática que busca resolver.}
  \criterion{La problemática es relevante y se relaciona con una necesidad real.}
  \criterion{El grupo identifica a quiénes afecta la problemática y cómo.}

  \category{2. Claridad y organización de la información}
  \criterion{La información se presenta de manera clara, coherente y organizada.}
  \criterion{Las ideas principales se comprenden fácilmente y siguen una secuencia lógica.}
  \criterion{El grupo utiliza ejemplos, datos o evidencias que ayudan a comprender la propuesta.}

  \category{3. Dominio y conocimiento del tema}
  \criterion{Los integrantes demuestran conocimiento del tema y de la problemática.}
  \criterion{Las respuestas a las preguntas son pertinentes y están bien fundamentadas.}
  \criterion{El grupo explica con sus propias palabras las razones de su propuesta.}

  \category{4. Innovación y pertinencia de la solución}
  \criterion{La propuesta ofrece una solución creativa o novedosa a la problemática.}
  \criterion{La solución responde directamente a la necesidad identificada.}
  \criterion{El grupo explica qué diferencia su propuesta de otras posibles soluciones.}

  \category{5. Viabilidad e impacto}
  \criterion{La propuesta parece viable considerando los recursos y las condiciones necesarias.}
  \criterion{El grupo identifica posibles beneficiarios y los cambios positivos que podría generar.}
  \criterion{Se consideran posibles dificultades, limitaciones o aspectos por mejorar.}

  \category{6. Comunicación y uso de recursos}
  \criterion{Los integrantes se expresan con claridad, seguridad y respeto.}
  \criterion{Los recursos de apoyo (diapositivas, prototipos, modelos, videos u otros) son
    pertinentes y facilitan la comprensión.}
  \criterion{Los recursos visuales o tecnológicos se utilizan adecuadamente, sin sustituir
    la explicación del grupo.}

  \category{7. Gestión del tiempo y trabajo en equipo}
  \criterion{La presentación se desarrolla dentro del tiempo asignado.}
  \criterion{El tiempo se distribuye adecuadamente entre la explicación de la problemática,
    la solución y sus beneficios.}
  \criterion{Se evidencia participación equilibrada y coordinación entre los integrantes.}
\end{tabularx}\par

\vspace{4mm}
\textbf{Retroalimentación breve.} Explica cómo funcionaría la solución alternativa o por qué ayudaría.\par
\vspace{3pt}
\setlength{\fboxsep}{4pt}
\setlength{\fboxrule}{0.4pt}
\responseblock{1}\hfill\responseblock{2}\par
\vspace{5pt}
\responseblock{3}\hfill\responseblock{4}\par

\end{document}
